\section{Evaluation Metrics}
\label{sec:evaluation_metrics}

This section defines the metrics used to evaluate the Gen-MOR models. Because the goal of this project is not only to reconstruct objective vectors but also to generate dominance-preserving plans, the evaluation includes both standard machine learning metrics and the dominance-specific metric introduced in Section~\ref{subsec:dvr}. The standard metrics measure reconstruction and classification behavior, while Dominance Validity Rate (DVR) measures whether generated samples remain structurally valid under the Pareto dominance relation.

\subsection{Reconstruction Mean Squared Error}
\label{subsec:metric_mse}

Reconstruction quality is measured using mean squared error (MSE). For each test sample $\mathbf{x}_i$, the model produces a reconstructed vector $\hat{\mathbf{x}}_i$. The reconstruction MSE is defined as
\begin{equation}
    \operatorname{MSE}
    =
    \frac{1}{N_{test}}
    \sum_{i=1}^{N_{test}}
    \left\|
        \mathbf{x}_i - \hat{\mathbf{x}}_i
    \right\|_2^2,
    \label{eq:metric_mse}
\end{equation}
where $N_{test}$ is the number of test samples. A lower MSE indicates that the model reconstructs the original objective vectors more accurately.

In Gen-MOR, MSE is useful because it shows whether the AE, VAE, and CVAE can learn the numerical structure of the multi-objective plan vectors. However, MSE alone does not measure whether generated samples preserve dominance relationships.

\subsection{Classification Accuracy}
\label{subsec:metric_accuracy}

Classification accuracy measures how often the model correctly predicts whether a plan belongs to $P^{+}$ or $P^{-}$. Let $y_i$ be the true label and let $\bar{y}_i$ be the predicted class label after thresholding the model output as defined in Eq.~\eqref{eq:classification_threshold}. Accuracy is defined as
\begin{equation}
    \operatorname{Accuracy}
    =
    \frac{1}{N_{test}}
    \sum_{i=1}^{N_{test}}
    \mathbb{I}
    \left[
        \bar{y}_i = y_i
    \right],
    \label{eq:metric_accuracy}
\end{equation}
where $\mathbb{I}[\cdot]$ is the indicator function.

Accuracy is easy to interpret, especially because the final dataset is near-balanced between $P^{+}$ and $P^{-}$. However, accuracy is still reported together with precision, recall, and F1 score because classification behavior may differ between the nondominated and dominated classes.

\subsection{Precision, Recall, and F1 Score}
\label{subsec:metric_precision_recall_f1}

Precision, recall, and F1 score provide a more detailed view of classification performance. In this experiment, the positive class corresponds to $P^{+}$, meaning that $y_i=1$ indicates a nondominated plan.

Precision measures the fraction of predicted positive samples that are actually positive:
\begin{equation}
    \operatorname{Precision}
    =
    \frac{TP}{TP+FP},
    \label{eq:metric_precision}
\end{equation}
where $TP$ is the number of true positives and $FP$ is the number of false positives.

Recall measures the fraction of actual positive samples that are correctly predicted:
\begin{equation}
    \operatorname{Recall}
    =
    \frac{TP}{TP+FN},
    \label{eq:metric_recall}
\end{equation}
where $FN$ is the number of false negatives.

The F1 score is the harmonic mean of precision and recall:
\begin{equation}
    \operatorname{F1}
    =
    2
    \cdot
    \frac{
        \operatorname{Precision}
        \cdot
        \operatorname{Recall}
    }{
        \operatorname{Precision}
        +
        \operatorname{Recall}
    }.
    \label{eq:metric_f1}
\end{equation}

F1 score is useful because it balances false positives and false negatives. In the Gen-MOR setting, this matters because confusing $P^{+}$ and $P^{-}$ is not merely a classification error; it also indicates that the model may not be learning the dominance boundary clearly.

\subsection{KL Divergence for VAE and CVAE}
\label{subsec:metric_kl_divergence}

For the VAE and CVAE models, the KL-divergence term measures how close the learned latent posterior is to the prior distribution. As defined in Eq.~\eqref{eq:vae_kl_loss} and Eq.~\eqref{eq:cvae_kl_loss}, the KL term regularizes the latent space toward the standard normal prior $p(\mathbf{z})=\mathcal{N}(\mathbf{0},\mathbf{I})$.

For a diagonal Gaussian posterior, the KL divergence is computed using the closed form from Eq.~\eqref{eq:vae_kl_closed_form}. A very small KL value may indicate that the latent space is close to the prior, but it may also suggest posterior collapse if the decoder ignores the latent variable. A very large KL value may indicate that the latent distribution is poorly regularized. Therefore, KL divergence is monitored together with reconstruction and classification losses during training.

\subsection{Dominance Validity Rate}
\label{subsec:metric_dvr}

Dominance Validity Rate is the main structural metric for Gen-MOR. It directly evaluates whether generated samples preserve the intended dominance behavior. The positive DVR, defined in Eq.~\eqref{eq:positive_dvr}, measures the fraction of generated $P^{+}$ samples that remain nondominated. The negative DVR, defined in Eq.~\eqref{eq:negative_dvr}, measures the fraction of generated $P^{-}$ samples that remain dominated.

The overall DVR is computed using Eq.~\eqref{eq:overall_dvr}. In the experimental results, class-specific DVR values are reported separately because positive and negative generation may have different difficulty levels. In particular, generating valid $P^{-}$ samples can be difficult because the generated sample must be dominated by at least one other plan under the selected objective subset.

\subsection{Multi-Seed Variance}
\label{subsec:metric_multiseed_variance}

To evaluate robustness, each model is trained and evaluated across multiple random seeds. For each metric, the mean and standard deviation are reported using Eq.~\eqref{eq:metric_mean} and Eq.~\eqref{eq:metric_std}. The mean indicates average performance, while the standard deviation indicates stability across runs.

This is especially important for generative modeling because different random seeds may produce different latent spaces and different generated samples. A model with high average DVR but high variance may not be reliable. Therefore, Gen-MOR reports both average performance and variability for reconstruction, classification, and dominance validity.

\subsection{Why Standard Metrics Are Insufficient}
\label{subsec:why_standard_metrics_insufficient}

Standard metrics such as MSE, accuracy, and F1 score are necessary but not sufficient for evaluating Gen-MOR. A model may reconstruct objective vectors accurately and classify test samples correctly, but still generate new samples that violate the dominance relation from Eq.~\eqref{eq:pareto_dominance}. This happens because dominance validity is relational. It depends on how a generated plan compares with other plans, not only on the generated vector itself.

For example, a generated vector may be close to the training distribution and may be classified as $P^{+}$ by the neural classifier. However, if another plan dominates it under the selected objective subset $S$, then it is not a valid generated $P^{+}$ sample. Similarly, a generated $P^{-}$ sample is invalid if no other plan dominates it. DVR is therefore required to measure the structural correctness of generated plans.